\chapter{Candidate bytes and committed values}
\MARSChapterHeading{Candidate bytes and committed values}
\label{revision-lab:encodings}

Use the complete \texttt{delta-codec.asm} and independent \texttt{delta-scan.asm}
in \texttt{revision/labs/files}. Read the main chapter's mapping, seven-bit groups
and prefix invariant first. These activities use ordinary Run and Debug; a native
guided lesson for this unit remains unfinished. Keep disposable copies for input
changes and intentional faults. Treat byte lengths, reading counts and capacities
as quantities with different units.

\section{Recover the wire from the readings}

Begin with the unchanged five-reading caller. Compute each exact difference from
its predecessor, map it to a nonnegative integer, then split it into base-128
digits. Set a continuation bit only when another digit follows. Explain why
1000 produces \texttt{d0 0f}, while a repeated reading produces \texttt{00}.
The entire wire is \texttt{d0 0f 04 01 00 ee a8 02}.

Break at \texttt{codec\_done}, before the caller publishes the results.
Live registers contain five consumed readings, eight produced bytes and status
zero in \texttt{\$v0/\$v1/\$a3}. Inspect the byte arena \texttt{stream}; the
three result words are still $-1$. At \texttt{results\_ready}, those words
contain five, eight and zero. Exit subsequently uses \texttt{\$v0}, so retain
the function result at the intended boundary.

To decode that stream in a fresh copy, set \texttt{method} to one, initialize
the eight bytes explicitly, set \texttt{byte\_count} to eight and retain word
capacity 32. Keep 192 total stream bytes by reducing its trailing filler to 184.
Bytes above 127 can be written as signed initializer spellings: 208 as $-48$,
238 as $-18$, and 168 as $-88$. This changes assembler source notation, not the
stored bit patterns. The decoder uses unsigned byte loads for interpretation.

At return, \texttt{\$v0}=8 bytes and \texttt{\$v1}=5 readings. Verify all five
words and the specified wire bytes. Count units reverse relative to the encoder.
Two routines can agree with each other while implementing the wrong format.

\section{Observe unused capacity after a complete prefix}

In an encoder copy, change only \texttt{byte\_capacity} from 160 to seven.
Derive each candidate length and the remaining capacity before running.
The first four readings need five bytes in total. The fifth needs three;
it cannot fit in the two remaining positions.

At \texttt{codec\_done}, expect four readings, five bytes and status one.
The first five stream bytes are \texttt{d0 0f 04 01 00}; every later byte remains
its initial hexadecimal \texttt{a5}. In particular, offsets five and six are
inside the offered capacity but outside the committed prefix. They must not
contain a partly emitted final value.

\begin{figure}[htbp]
\centering
\begin{tikzpicture}[alt={Actual light-theme Memory row at stream address hexadecimal 100100c0 after encoding with byte capacity seven. Bytes zero through four are d0 0f 04 01 00. The following eleven displayed bytes remain a5. A left native box marks the five committed bytes; a right box marks offsets five and six, which are offered capacity but cannot fit the rejected three-byte candidate.}]
\node[anchor=south west,inner sep=0] (a) {\MARSRevisionCrop{revision/captures/delta/delta-atomic-prefix.png}{0.97\linewidth}{0.0328125}{0.3388889}{0.584375}{0.6088889}};
\begin{scope}[shift={(a.south west)},x={($(a.south east)-(a.south west)$)},y={($(a.north west)-(a.south west)$)}]
\draw[black,line width=.7pt] (.208,.02) rectangle (.436,.43);
\draw[black,line width=.7pt] (.445,.02) rectangle (.532,.43);
\end{scope}
\end{tikzpicture}
\caption{Actual Debug memory after the capacity-seven encoder returns. The left box marks the committed prefix; the right marks two unused positions within capacity. Neither those positions nor later backing storage may contain a rejected candidate's temporary output. Native LaTeX supplies the crop and boxes.}
\label{fig:revision-delta-native-atomic}
\end{figure}


Inspect \texttt{enc\_check\_capacity} on a fresh execution. Its candidate length
is three and remaining capacity is two. The function chooses \texttt{enc\_full}
before computing and using the candidate's output cursor. Earlier accepted
readings remain committed. A final-memory comparison is not sufficient to prove
that a rejected candidate was never written; the public checks also observe
memory accesses, including stores later restored.

Repeat with capacities zero, one, two, five and eight. For each, predict both
return counters, status and the first unchanged byte. Then set reading count
zero and byte capacity zero together. Empty input succeeds without reading data;
nonempty input with zero capacity reports capacity before reading its first value.
Explain the ordering from the loop's top tests.

\section{Distinguish scanning from committing}

Prepare a decoder copy with method one, byte count three and stream
\texttt{02 80 00}, followed by 189 filler bytes. Keep the default word arena
so an unmodified second word is recognizable as 1002. The first byte maps to
difference one and commits reading one. The next candidate represents unsigned
zero using two bytes, which violates shortest spelling.

\Needspace{6\baselineskip}
Break at \texttt{dec\_bad}, before its status instruction. General register
\texttt{\$t0} is scan index three, and \texttt{\$t1} is candidate unsigned
zero. In the same stopped state, predecessor \texttt{\$s4}, committed input
\texttt{\$s5} and committed readings \texttt{\$s6} are all one. The first
word has become one; the second remains 1002. These are separate observations
of a rejected candidate and the accepted prefix, not conflicting counter values.

\begin{figure}[htbp]
  \centering
  \begin{tikzpicture}[alt={Two actual debugger views of the same dec_bad stop for bytes 02 80 00. The candidate scan registers show t0 equal to three and t1 equal to zero. The committed state shows s4, s5 and s6 each equal to one: last valid reading, committed bytes and committed readings. Native boxes distinguish the candidate from the accepted prefix. The rejected redundant spelling has been examined but has not been committed.}]
    \node[anchor=north west,inner sep=0,font=\small\bfseries] (atitle) at (0,0) {1. Candidate scan};
    \node[anchor=north west,inner sep=0] (a) at ([yshift=-2mm]atitle.south west) {\MARSRevisionCrop{revision/captures/delta/delta-candidate.png}{0.29\linewidth}{0.24609375}{0.180}{0.66015625}{0.77555556}};
    \node[anchor=north west,inner sep=0,font=\small\bfseries] (btitle) at (6,0) {2. Committed prefix};
    \node[anchor=north west,inner sep=0] (b) at ([yshift=-2mm]btitle.south west) {\MARSRevisionCrop{revision/captures/delta/delta-committed.png}{0.29\linewidth}{0.24609375}{0.0788889}{0.66015625}{0.8544444}};
    \begin{scope}[shift={(a.south west)},x={($(a.south east)-(a.south west)$)},y={($(a.north west)-(a.south west)$)}]
      \draw[black,line width=.7pt] (.035,.02) rectangle (.50,.98);
    \end{scope}
    \begin{scope}[shift={(b.south west)},x={($(b.south east)-(b.south west)$)},y={($(b.north west)-(b.south west)$)}]
      \draw[black,line width=.7pt] (.035,.02) rectangle (.50,.98);
    \end{scope}
  \end{tikzpicture}
  \caption{Actual decimal register observations before the malformed-status instruction and before saved-register restoration. Scan index three includes the rejected two-byte candidate; committed input remains one. These views come from the same source and stop in fresh deterministic executions. Native LaTeX supplies labels, crops and boxes.}
  \label{fig:revision-delta-native-candidate}
\end{figure}


Continue to \texttt{codec\_done}. The public return is one committed byte,
one reading and status two. The routine has restored saved registers, so its
working values in \texttt{\$s4/\$s5/\$s6} no longer describe the loop state.
Use the return registers after this boundary. Advancing committed input to the
scan index would incorrectly claim the rejected candidate had been consumed.

Now reduce word capacity to one. The function accepts the first byte and stops
with status one, without examining either malformed tail byte. This is a
capacity result, not a malformed result, because the next candidate is outside
the amount of output the caller allowed the function to attempt.

\section{Separate byte syntax from reading range}

Classify these streams before executing them from predecessor zero:
\begin{verbatim}
00                 zero, one byte
80 01              unsigned 128, difference 64
80 00              redundant spelling of unsigned zero
80                 continued value cut off by input end
80 80 80 80 20     fifth payload too wide
80 80 80 80 80     fifth byte still requests continuation
ff ff ff ff 1f     valid spelling, reconstructed reading out of range
\end{verbatim}
The last three examples distinguish a width rule, a termination rule and a
signed reading rule. For a malformed fifth continuation, the decoder must not
read a sixth byte even if padding is mapped. For \texttt{ff ff ff ff 1f},
unsigned value 8589934591 reverses to difference $-4294967296$; only the reading
range check rejects it. Use status two for byte syntax and three for range.

Next encode the readings minimum word, maximum word, minimum word. The second
difference is 4294967295 and the third is $-4294967295$. Their mapped spellings
are \texttt{fe ff ff ff 1f} and \texttt{fd ff ff ff 1f}. The first reading
uses \texttt{ff ff ff ff 0f}. Explain why a word subtraction can pass small
examples while corrupting these transitions. Find the first incorrect arithmetic
result before investigating the later byte loop.

The kernel, caller and same-session checks accept equivalent complete codecs:
\Needspace{15\baselineskip}
\begin{verbatim}
java -jar /path/to/cli.jar grade \
  packages/023-delta-varint/checks/kernel.kts \
  revision/labs/files/delta-codec.asm
java -jar /path/to/cli.jar grade \
  packages/023-delta-varint/checks/program.kts \
  revision/labs/files/delta-codec.asm
java -jar /path/to/cli.jar grade \
  packages/023-delta-varint/checks/reuse.kts \
  revision/labs/files/delta-codec.asm
\end{verbatim}
The arena sizes do not enlarge the selected slices. Preserve full pointers,
all saved register bits and stack state. The worked functions accept the
original package's eight-byte-aligned incoming stack; the ordinary caller uses
sixteen-byte alignment. A new call begins from predecessor zero even after a
partial or failed call.

\section{Compare bytes separately from execution work}

Run the repeatable public measurement tool with a new output directory:
\begin{verbatim}
python3 tools/measure_revision_delta.py \
  --java /path/to/jdk-25/bin/java --cli /path/to/cli.jar \
  --output /tmp/my-delta-measurements
\end{verbatim}
It refuses an existing directory, checks the complete encoding and decodes its
actual bytes in the same execution session. Its receipt binds source, public
oracle, tool and CLI hashes. It records representation bytes without framing
and completed instructions for each function's entry through return, including
frame operations and validation. Assembly, session startup, caller publication,
host inspection and hardware cycles are outside that interval.

Explain why the steady dataset benefits from deltas, the small ramp gains
nothing over independent mapping, and the spike loses two bytes. Do not compare
the encoder instruction count with a raw word copy that has different work and
then call the ratio a machine speedup. If investigating another implementation,
keep numerical semantics, input, capacity, validation and measurement interval
consistent. A correct shorter wire and a correct faster implementation are
separate findings.

\section{Build a read-only scanner independently}

Implement \texttt{scan\_delta} in \texttt{delta-scan.asm}, retaining its caller
and data harness. Inputs are full byte pointer \texttt{\$a0}, byte length zero
through 160 in \texttt{\$a1}, and reading limit zero through 32 in
\texttt{\$a2}. Outputs are committed byte and reading counts in
\texttt{\$v0/\$v1}, last committed signed reading in \texttt{\$a2}, and
status in \texttt{\$a3}. With no committed reading, the last value is zero.

Apply the decoder's exact mapping, shortest-spelling, fifth-byte, signed-range
and precedence rules. Initialize predecessor on every call. Read only the
selected input slice, and no caller bytes when input is empty or limit is zero.
Write no caller bytes, including temporary writes restored before return.
Private stack storage and submission-owned constants are permitted. Preserve
saved registers, frame pointer and stack pointer, return through the supplied
continuation and perform no console I/O.

The unchanged stream returns eight bytes, five readings, last reading 20000 and
status zero at \texttt{scan\_ready}. The caller publishes them only afterward.
Choose a valid prefix whose last value differs from the rejected candidate;
use it to check that rejection retains the earlier last reading. Also repeat a
successful call after an error to expose accidental retained state.
\Needspace{6\baselineskip}
Run the public checks:
\begin{verbatim}
java -jar /path/to/cli.jar grade \
  revision/labs/checks/delta-scan.kts \
  /path/to/your/delta-scan.asm
\end{verbatim}
Keep your candidate/prefix invariant, one rejected-candidate observation and an
explanation of error precedence with the implementation. Behavioral checks do
not require the worked decoder's temporary registers or source spelling.
